\documentclass[11pt]{article}
\usepackage[margin=1in]{geometry}
\usepackage{amsmath,amssymb,amsthm,booktabs,array}
\usepackage[hidelinks]{hyperref}
\usepackage{microtype}
\newtheorem{proposition}{Proposition}
\title{The Value of Fleet-Resource Coupling: Objective Levels and Physical Realization}
\author{A research section accompanying the stable GI-TBC method}
\date{}
\begin{document}
\maketitle

\section{Inventory proof and the physical optimum}
The stable Gini-interval tailored branch-and-cut method (GI-TBC) maintains
complete improving-domain coverage, a ledger of scope-valid bounds, and
independently verified fleet incumbents. Integer-inventory pricing can
strengthen the ledger without changing the interval controller or route
producer. This separation makes it possible to study a stronger relaxation
without silently making it part of the complete algorithm. The default
method remains GI-TBC/B; E+ is the enhanced comparison baseline inherited
from the original-objective inventory study.

Let $r_i=c_iY_i$, $S=\sum_i r_i$, and
\[
 G(Y)=\frac{\sum_{i<j}|r_i-r_j|}{nS},\qquad
 H(Y)=\sum_i\omega_i|r_i-1|,\qquad J(Y)=G(Y)+H(Y).
\]
The original zero convention is retained. The coefficient $\omega_i$
already includes $\lambda$. All domains below lie in the same qualified
improving inventory scope, with the original supports, hard total and
coverage rules. The present scopes have $\gamma_L=0$ and a strictly
positive stock denominator.

Write $D_I$ for the integer inventory relaxation. For RX35, let $D_F$ delete
the 26 vectors excluded by the two preceding RX35 family certificates.
This deletion set is specific to RX35; for the NB35 comparison take
$D_F=D_I$. Later family exclusions do not silently redefine $D_F$.
Let $D_R(T)$ be the inventory projection of free vehicle assignment,
per-vehicle pickup/drop and return-load balance, handling, and valid
closed-tour lower bounds for at most three stations. Finally, $D_P(T)$
contains inventories realized by a complete original fleet. Under the
inherited physical and numerical qualification,
\[
 D_P(T)\subseteq D_R(T)\subseteq D_F\subseteq D_I,
 \qquad F_I\le F_F\le F_R\le F_P,
\]
where $F_D=\min_{Y\in D}J(Y)$. A strict set inclusion does not by itself
make the corresponding minimum strictly larger. In particular, the best
saved inventory families are not a proved ordering of the global
$k$-best inventories.

\begin{proposition}[A resource witness bounds same-source cut value]
Suppose $L_I\le F_I$ and a fully checked original-$T$ resource assignment
realizes inventory $Y^R$ of price $U_R=J(Y^R)$. If a cut domain $D_C$
satisfies $D_R(T)\subseteq D_C\subseteq D_I$, then
\[
 F_C\le F_R\le U_R,\qquad
 0\le F_R-F_I\le U_R-L_I.
\]
\end{proposition}
\begin{proof}
Domain containment orders the minima; $Y^R$ supplies an actual member of
the resource domain. The inventory lower bounds $F_I$ from below.
\end{proof}

Thus a low resource witness limits every original-objective floor derived
solely from rearranging the same resource facts, including projection,
core and no-good cuts. The limit does not apply to stronger travel or
prefix-load information, or to producing a real route. An inventory-only
upper cannot play the role of $U_R$. A known physical route supplies a
valid but potentially loose resource upper by projection; a resource
assignment alone supplies no physical upper.

\section{Whole objective levels, rather than isolated exclusions}
The principal question can be posed as the hard level
\[
 \mathcal R(\tau,T)=\{(Y,a):Y\in D_I,\ J(Y)\le\tau,
               \ a\text{ satisfies the full necessary resource model}\}.
\]
A qualified infeasibility result excludes the complete level, so
$F_R>\tau$ and $F_P>\tau$. Within the inherited proof policy, a conservative
recorded lower is $\tau$. Feasibility instead requires a complete saved
integer inventory and assignment, followed by independent evaluation of
the true objective and original time conditions. An unfinished query
does not improve either side of the prior interval.

The predeclared primary level is
\[
 \tau_* = L_I+\tfrac14(U_P-L_I).
\]
For RX35 this is approximately $0.2519564273646922$. Excluding it would
remove one quarter of the inherited physical gap. Since the known
inventory upper is approximately $0.23612171136420015$, that conclusion
would also establish a structural resource/inventory difference of about
$0.015834716$. These are implications of a complete-level proof, not
consequences of excluding further stored points.

The normalized inventory formulation uses a Gini epigraph $q$. At an
integer stock with $\gamma_L=0$, $q=G(Y)$ is available; hence
$q+H\le\tau$ covers every true low-objective stock. For each station choose
an exact coefficient anchor $a_i$ and a power-of-two scale $s$. Actual
emitted coefficients $\widehat c_{iy}$ represent
$(h_i(y)-a_i)/s$, rounded downward except for deliberate tiny omissions.
Let
\[
 \epsilon=\sum_i\max\bigl(0,\max_y\{
 \widehat c_{iy}-(h_i(y)-a_i)/s\}\bigr).
\]
The actual cutoff is
\[
 q/s+\sum_{i,y}\widehat c_{iy}x_{iy}
 \le \operatorname{up}\!\left((\tau-\sum_i a_i)/s+\epsilon\right).
\]
Every coefficient and the RHS are checked in the Native matrix. The
one-hot relations bound the total positive coefficient excess, so no
omitted tiny coefficient silently deletes a valid low-objective point.
A returned point is still repriced exactly; a relaxed cutoff value is
not sufficient to assert membership in the true level. The old
known-lower auxiliary row is not used in these level models.

\section{A qualified integer budget expression}
The resource projection inherits classical vehicle-resource and tour
necessary conditions. Its integer rounding follows the earlier rounded
handling and anchor/whole-handling studies; it is not a new general cut
theory. The old production H-ACT mechanism remains parked.

Let $h>0$ be a common lower handling coefficient and $\ell(S)$ a valid
closed-tour lower for visiting $S$. Raw physical and emitted canonical
coefficients first share a coefficient-wise lower envelope. The
shortest-path and subset-tour computations use exact binary-rational
arithmetic and downward enclosures. A common upper horizon $\overline T$
covers the original physical acceptance rule, its arithmetic envelope,
and the inherited time guard. Both tested expressions use the identical
coefficient package; a bare floor at nominal $T$ is not substituted.

Total pickups $P_k$ are integral sums of integer quantities and binary
service decisions. Define
\[
 U=\left\lfloor\frac{\overline T}{h}\right\rfloor,\qquad
 b_S=\left\lfloor\frac{\overline T-\ell(S)}h\right\rfloor.
\]
Always $P_k\le U$. For $0\le b_S<U$, replace the real-valued time row by
\[
 P_k+(U-b_S)\sum_{i\in S}s_{ki}
 \le b_S+(U-b_S)|S|.                         \tag{1}
\]
If $b_S<0$, impose $\sum_{i\in S}s_{ki}\le |S|-1$.
Total pickups are never incorrectly bounded by $Q_k$: deliveries can
free space for later pickups, while the physical capacity applies to
prefix loads and the returned load.

\begin{proposition}[Integer equivalence and subset dominance]
In the same coefficient package, (1), the forbidden-set rows and
$P_k\le U$ preserve the integer service domain and imply the old
real-valued resource rows in the LP relaxation. If a proper subset
$R\subset S$ has $b_R\le b_S$, its budget row and the global budget
imply the $S$ row. A forbidden subset forbids its supersets.
\end{proposition}
\begin{proof}
Put $\delta_S=|S|-\sum_{i\in S}s_{ki}$. At integer visits, $\delta_S=0$
activates the exact integer budget, whereas $\delta_S\ge1$ leaves the
global budget sufficient. For $0\le\delta_S\le1$, the new bound connects
the downward-rounded budgets at both endpoints, and therefore lies below
the old real-valued time bound. For $\delta_S\ge1$, $P_k\le U$ already
implies that old bound. For dominance, $\delta_R\le\delta_S$ and, on
$[0,1]$, $b_R+(U-b_R)\delta_S\le b_S+(U-b_S)\delta_S$; the global budget
covers the remaining range. The forbidden-set statement follows directly.
\end{proof}

Under the qualified acceptance horizon, RX35 retains 4087 of 7175 subset
rows per vehicle; its 264 negative budgets reduce to 81 minimal forbidden
sets. NB35 retains 4755 subsets, with four minimal forbidden sets. Two
RX35 and five NB35 budgets cross an integer threshold when the time
acceptance margin is included. These counts explain the scale of the
expression opportunity; they establish no speedup on their own.

\subsection{A fractional ownership obstruction}
An integer stock can be difficult for the resource model even when its
inventory variables already have an exact normalized lift. The following
elementary construction separates inventory integrality from vehicle
ownership integrality. It concerns the explicit bare LPs, before additional
integer-valid cuts or branching.

\begin{proposition}
Let an integer stock have an inventory LP lift in the hard level. Suppose
there are $M\ge2$ vehicles, every nonzero service is individually eligible
for every vehicle, and its total pickups/drops $P,D$ satisfy
\[
 0\le(P-D)/M\le\min_k Q_k,\qquad P\le(M-1)U.
\]
If every singleton tour lower is at most $\overline T$, assigning a fraction
$1/M$ of each nonzero service to every vehicle extends that stock to a
feasible LP point of both ideal resource expressions, including all their
rounded and forbidden-subset rows.
\end{proposition}
\begin{proof}
Each vehicle receives pickup $P/M$, drop $D/M$ and visit value at most
$1/M$ per station. Assignment and return balance hold. The global budget
holds since $P/M\le(M-1)U/M\le U$. For a singleton, the old time row is at
most $(hP+\overline T)/M\le\overline T$. For $|S|\ge2$,
$\sum_{i\in S}s_{ki}-|S|+1\le0$, so handling alone suffices for the old
row. For a rounded nonforbidden singleton,
$P/M\le(M-1)U/M\le b_S+(U-b_S)(1-1/M)$.
For larger sets, $\delta_S\ge1$ makes the global budget sufficient.
A forbidden set has size at least two under the singleton assumption,
and $\sum_{i\in S}s_{ki}\le |S|/M\le |S|-1$.
\end{proof}

This is a standard fractional-assignment weakness, not an integer resource
upper. A checked instantiation at price $J(Y)$ can upper-bound the
original-price optimum of those fixed bare LPs. It cannot upper-bound
$F_R$, and additional integer-valid cuts can remove it. In particular,
survival of one impossible stock in the LP does not establish a strict
gap between the global inventory and resource minima. Actual binary-matrix
coefficients and RHS values require a separate saved-point check before
an exact numerical LP witness is reported.

\section{Finite queries and a complete hard-level campaign}
Two preselected free-assignment queries were tested first. RX35's retained
Y3 exact-price family contains ten members; NB35's saved-best family is a
singleton. Both were excluded at their original/guarded time horizon.
The queried v1 distance reconstruction gives necessary-time lowers
$6310.129749041621$ and $7206.982290055626$ seconds, respectively, against
original horizons $6300$ and $7200$. These numeric time lowers belong to
the queried v1 models. A finite audit of the actual C++ physical duration
rows found identical integer budgets at both nominal and guarded horizons,
which transfers the fixed-time exclusions; it does not relabel the old
numbers as identical unconditional C++-coefficient time minima. Their
relaxed-time assignments are neither original-time resource uppers nor
physical routes. Family scope remains explicit.

\begin{table}[htbp]
\centering
\begin{tabular}{lrrr}
\toprule
Role & $L_I$ & $U_I$ & Inherited $U_R=U_P$\\
\midrule
RX35 & .2360521029 & .2361217114 & .2996694007\\
NB35 & .3959564087 & .3960465490 & .4502915033\\
\bottomrule
\end{tabular}
\caption{Retained coarse root-scope intervals. The physical upper projects
to a resource upper; no novel low-price integer resource assignment or
new physical upper was obtained. RX35's $D_F$ upper remains the
inventory-only Y3 price, approximately $.2368540093$.}
\end{table}

The one qualified expression alternative retains $4087$ RX35 subsets
instead of $7175$, reducing the full model from $24154$ to $14890$ rows
while keeping $2673$ columns. A finite raw diagnostic returned at $594$
nodes in $753.307$ outer seconds; the rounded diagnostic reached $512$
nodes in $847.824$ seconds. These are unequal observed node counts, not a
matched-node speed comparison. No compression speedup was demonstrated,
and the raw expression was selected for the substantive capacity.

The actual C++-coefficient raw hard level at
$\tau_*=.2519564273646922$ then completed its full declared $129600$-second
Native window. All inventory and vehicle decisions remained free, with
Seed $0$, one thread, MIPFocus $1$ and SolutionLimit $1$. The objective was
constant-zero feasibility under the hard original-price cutoff. Original
model $T=6300$ seconds is separate from the experimental window.
The call returned normally with status TIME\_LIMIT, no integer solution,
$258972$ final nodes, $160559495$ simplex iterations and
$318526.526959$ work units. Native Runtime was $129600.004$ seconds and
the charged outer time was $129601.832$ seconds. Its zero feasibility bound
is not an original-$J$ lower. The full level is therefore UNKNOWN and the
table's prior intervals are retained.

After that call closed, an exact rational saved-point check instantiated
the fractional-ownership construction at Y3. Its total pickups and drops
are both $56$; each vehicle owns $1/3$ of every changed station and has
zero returned load. The same point satisfies every actual row and bound
of both admitted matrices, with exact maximum violation zero. It has
$102$ fractional integer columns, maximum integrality violation $1/3$,
and original price approximately $.23685400930608236$. A separate saved
readback repeats all exact row and bound checks. Thus the two bare LPs
retain this integer stock through fractional ownership even though its
integer resource family is excluded. This is no integer resource upper,
no limit on stronger integer-valid cuts, and no proof that the global
resource minimum strictly exceeds the inventory minimum. It documents a
specific obstruction without identifying the unique cause of search cost.

The total scientific outer charge rounds to $37.1$ hours, including failed
preparation, the normally stopped watchdog/coefficient-repair calls,
necessary qualification, and saved-original readers. All receipts are
closed; nested Runtime and Work are not charged again. The first exact
LP reader failed on an infinite-bound conversion; its original bytes and
fee remain, and only the signed infinite-bound readback was corrected.
PDF compilation/rendering, pure writing and publication I/O are separate.
Current results and exact final costs accompany the saved evidence.

This completes a substantive finite computational negative under the
predeclared capacity. Whether $F_R-F_I$ is worthwhile remains open.
No new original-objective lower or true route benefit justifies freezing
E+R, so the conditional eight-arm panel was not entered. The next research
question concerns justified information about integer vehicle ownership
or stronger physical travel/load structure. Stable GI-TBC/B remains the
default and E+ the enhanced baseline; no production change or complete-arm
benefit is claimed.


\section{Evidence consumption and the complete method}
A whole-level resource lower is a bound for the original inventory scope
and can enter the GI-TBC ledger only with that scope and its coverage
qualification. A local family proof cannot close all low-objective
inventories. A physical route can update $U$ only after the unchanged
original-time, station, vehicle, operation and prefix-load checks. A
relaxed-time assignment, a log objective or an unfinished callback point
does not substitute for either certificate.

Any complete E+R candidate must pay for its own inventory pricing,
resource-level computation, route construction and verification. Offline
stocks, old assignments and proofs are not free starts for a nominal
fresh arm. A uniform rule may use current mathematical bounds and fixed
node/iteration capacities; it may not dispatch on instance names,
historical wins or a timer switch. A worthwhile structural finding need
not become a production mechanism. Stable GI-TBC exposition and the
maintained paper continue independently of whether this research branch
produces such a candidate.

\paragraph{Attribution and scope.}
This section continues the support-duration studies (R51--52), the
free-fleet/time-oracle studies (R61), the
physical-envelope rounded handling repair (R92), full handling with
anchors (R103), resource propagation (R150--151), and integer-inventory
objective proof consumption (R152--157). Results, scope identities,
actual Native matrices, source snapshots, fees and same-agent separate
readbacks accompany the present finite study. No reviewer independence,
universal performance benefit or default promotion is inferred.
\end{document}
